% Zdroj: PDF priklady-leto-2025.pdf, fyzická str. 22-23. Značka: specializace UI-SU, UI-ZPJ, UI-ROB. Zařazeno: S1 (MDP / Bellmanova rovnice).
\section{Bellmanova rovnice užitku}
\szzotazka{léto 2025}{29}{Bellmanova rovnice užitku (MDP, iterace strategií)}

\noindent\textbf{Zadání.}\par
Robot se pohybuje v konečném stavovém prostoru $S$. Pro každý stav $s \in S$ známe množinu možných akcí $A(s)$ a pravděpodobnost $P(s' \mid s, a)$, že se robot ze stavu $s$ akcí $a \in A(s)$ přesune do stavu $s' \in S$. Při průchodu stavem $s \in S$ dostane robot odměnu (reward) $R(s)$. Celkový užitek ze stavu $s'$ je diskontován faktorem $\gamma = 0.9$.

\begin{enumerate}[label=(\arabic*),nosep]
  \item Napište rovnici pro výpočet maximálního užitku (utility) $U(s)$ ve stavu $s$.
\end{enumerate}

\begin{center}\includegraphics[width=0.8\linewidth]{mdp-stavovy-graf.png}\end{center}

Obrázek udává stavový prostor, odměnu v jednotlivých stavech a dvě možné akce z každého stavu kromě koncových stavů $E$ a $F$. Pravděpodobnost, že se robot přesune do stavu, který je daný orientovanou hranou zvolené akce, je 0.8 a ve zbylé pravděpodobnosti 0.2 se přesune do stavu, do kterého směřuje druhá možná akce.

\begin{enumerate}[label=(\arabic*),nosep,start=2]
  \item Napište soustavu rovnic, jejichž řešení dá užitky jednotlivých stavů, jestliže zadáme robotu akce (strategii/policy) $B \to C$, $C \to D$ a $D \to F$.
  \item Jakým algoritmem bychom tuto soustavu vyřešili?
  \item Jak ze spočtených užitků zjistíme, zda zadaná volba akcí dává maximální možné užitky pro všechny stavy?
  \item Jak bychom pokračovali, abychom našli akce pro všechny stavy maximalizující užitek?
\end{enumerate}

\medskip
\noindent\textbf{Řešení.} \textit{(V~PDF bez náčrtu řešení; vlastní vzorové řešení, soustava ověřena numericky.)}\par

\noindent Stavy a~odměny: $B(-3),C(1),D(1)$ (nekoncové), $E(5),F(6)$ (koncové). Akce: $B$ má $\{\to C,\to D\}$, $C$ má $\{\to D,\to E\}$, $D$ má $\{\to F,\to B\}$; zvolená akce vede do~svého cíle s~pravd.\ $0{,}8$, do~cíle druhé akce s~pravd.\ $0{,}2$. Koncové užitky $U(E)=R(E)=5$, $U(F)=R(F)=6$.

\noindent\emph{(1) Bellmanova rovnice} (maximální užitek):
\[ U(s)=R(s)+\gamma\,\max_{a\in A(s)}\sum_{s'} P(s'\mid s,a)\,U(s'). \]

\noindent\emph{(2) Soustava pro strategii} $B\to C,\ C\to D,\ D\to F$ (pevná strategie $\Rightarrow$ odpadá $\max$, rovnice jsou lineární):
\[
\begin{aligned}
U(B)&=-3+0{,}9\,[\,0{,}8\,U(C)+0{,}2\,U(D)\,],\\
U(C)&=\ \ 1+0{,}9\,[\,0{,}8\,U(D)+0{,}2\cdot5\,],\\
U(D)&=\ \ 1+0{,}9\,[\,0{,}8\cdot6+0{,}2\,U(B)\,].
\end{aligned}
\]

\noindent\emph{(3) Algoritmus řešení.} Jde o~\emph{vyhodnocení strategie} (policy evaluation): tři lineární rovnice o~třech neznámých vyřešíme přímo (Gaussova eliminace, $O(n^3)$), případně iterativně zjednodušeným Bellmanovým updatem bez~$\max$, dokud se užitky neustálí. Řešení:
\[ \boxed{U(B)\approx 2{,}383,\quad U(C)\approx 6{,}039,\quad U(D)\approx 5{,}749},\qquad U(E)=5,\ U(F)=6. \]

\noindent\emph{(4) Test optimality.} Pro každý stav ověříme \emph{Bellmanovo optimalitní kritérium}: zda zvolená akce dosahuje $\max_a \sum_{s'} P(s'\mid s,a)U(s')$, tj.\ zda by krok zlepšení strategie nezměnil akci. Spočítáme očekávaný užitek obou akcí (z~užitků výše): $B$: $\to C$ dá $5{,}981$ vs.\ $\to D$ $5{,}807$; $C$: $\to D$ dá $5{,}599$ vs.\ $\to E$ $5{,}150$; $D$: $\to F$ dá $5{,}277$ vs.\ $\to B$ $3{,}106$. V~každém stavu je zvolená akce lepší, takže strategie \textbf{je optimální}.

\noindent\emph{(5) Hledání optimálních akcí.} Obecně \emph{iterace strategií} (policy iteration): střídáme vyhodnocení strategie (bod~3) a~zlepšení strategie (v~každém stavu vezmi akci s~nejvyšším očekávaným užitkem, bod~4), dokud se strategie přestane měnit. Zde jsme v~bodě~4 ověřili, že daná strategie je optimální, takže iterace strategií by skončila hned.
